%%%%%%%%%%%%%%%%%%%%%%%%%%%%%%%%%%%%%%%%%%%%%%%%%%%
\section{Основные понятия}
%%%%%%%%%%%%%%%%%%%%%%%%%%%%
\subsection{Единицы измерения}
%%%%%%%%%%%%%%%%%%%%%%%%%%%%%%%%%%%%%%%%%%%%%%%%%%%%
%%%%%%%%%%%%%%%%%%%%%%%%%%%%%%%%%%%%%%%%%%%%%%%%%%%%
\z{Запишите в стандартном виде:\par
\hpu{$6400 \cdot 10^3$}
\hpu{$0.0025$}
\hpu{$602000 \cdot 10^{18}$}\par
\hpu{$0.00302 \cdot 10^{-5}$}  
\hpu{$2398 \cdot 10^{-2}$}
\hpu{$0.1309 \cdot 10^{12}$}
}
%
{\puan $6.4\cdot10^{6}$;~ \puan $2.5\cdot10^{-3}$;~ \puan $6.02\cdot10^{23}$;

 \puan $3.02\cdot10^{-8}$;~ \puan $2.398\cdot10^{1}$;~  \puan $1.309\cdot10^{11}$}

%%%%%%%%%%%%%%%%%%%%%%%%%%%%%%%%%%%%%%%%%%%%%%%%%%%%
\z{Переведите в систему СИ и выразите в стандартном виде:\par
\hpu{$6400 \cdot 10^3 \, \text{км}$}
\hpu{$1079252848.8 \, \text{км/ч}$}
\hpu{$54 \cdot 10^5 \, \text{мм}^3$ }\par
\hpu{$0.0034 \, \text{кг}$}  
\hpu{$24 \, \text{ч}$}
\hpu{$0.23 \, \text{км}^2$}
}
%
{\puan $6.4\cdot10^{9}\,\text{м}$;~ \puan $2.99792468\cdot10^{8}\,\text{м/с}$;~ \puan $5.4\cdot10^{-3}\,\text{м}^3$; 

\puan $3.4\cdot10^{-3}\,\text{кг}$;~ \puan $8.64\cdot10^{4}\,\text{с}$; ~\puan $2.3\cdot10^{5}\,\text{м}^2$}

%%%%%%%%%%%%%%%%%%%%%%%%%%%%%%%%%%%%%%%%%%%%%%%%%%%%
\z{Капля объёмом $0.01 \, \text{мм}^3$ растеклась по поверхности $S = 10 \, \text{см}^2$. Оцените размер молекулы масла.}
%
{$d=\frac{V}{S}=10^{-8}\,\text{м}=10\,\text{нм}$}


%%%%%%%%%%%%%%%%%%%%%%%%%%%%%%%%%%%%%%%%%%%%%%%%%%%%
\z{Зеркальная поверхность $S = 1 \, \text{м}^2$ покрыта серебром массой $0.1 \, \text{г}$. На зеркале $N = 100$ слоёв серебра. Оцените размер атома серебра.}
%
{$d \approx 10^{-10}\,\text{м}$}

%%%%%%%%%%%%%%%%%%%%%%%%%%%%%%%%%%%%%%%%%%%%%%%%%%%%
\z{На палубе судна имеется прямоугольная площадка размерами $10 \times 15 \, \text{м}$. Сколько контейнеров может поместиться на площадке, если контейнер представляет собой куб с длиной ребра $2 \, \text{м}$?}
%
{35 контейнер}

%%%%%%%%%%%%%%%%%%%%%%%%%%%%%%%%%%%%%%%%%%%%%%%%%%%%
\z{В древнем Риме было несколько единиц измерения объёма: секстарий, киафа, урна и амфора. Известно, что $12$ киафов составляют $1$ секстарий, $1$ урна равна половине амфоры, а в $1$ урне $288$ киафов. Определите, что больше по объёму: $3$ амфоры или $150$ секстариев.}
%
{$150$ секстариев больше, чем $3$ амфоры}

%%%%%%%%%%%%%%%%%%%%%%%%%%%%%%%%%%%%%%%%%%%%%%%%%%%%
\z{Одной из основных единиц земельной меры выбрали акр, который равен той площади, которую способен вспахать крестьянин на одном быке за один день. Английская единица площади акр --- равна $4046.86 \, \text{м}^2$. Сколько быков потребуется, чтобы вспахать за сезон (длительностью $15$ дней) поле площадью $2000$ соток, если $1$ сотка = $100 \, \text{м}^2$? Ответ округлить до целого в большую сторону. }
%
{требуется $4$ быка}

